\subsection{The Bicommutant of a Generating Module}

THEOREM 2. --- A generating module is balanced.

Let A be a ring, and let M be a generating A-module; by definition, there exist an integer \(n \geqslant 0\) , elements \(x_{1}, \ldots, x_{n}\) of M, and elements \(x_{1}^{*}, \ldots, x_{n}^{*}\) of the dual \(M^{*}\) of M satisfying \(\sum_{i=1}^{n} \langle x_{i}, x_{i}^{*} \rangle = 1\) . Recall (II, §4, No.~2, p.~271) that we define a group homomorphism \(\theta : M^{*} \otimes_{A} M \to \operatorname{End}_{A}(M)\) by the formula \(\theta(x^{*} \otimes y)(x) = \langle x, x^{*} \rangle y\) . If u is an element of the bicommutant of M, then it commutes with \(\operatorname{End}_{A}(M)\) ; therefore, for every \(y \in M\) , we have

\[
\begin{array}{l} u (y) = u \Big (\sum_ {i = 1} ^ {n} \langle x _ {i}, x _ {i} ^ {*} \rangle y \Big) = \sum_ {i = 1} ^ {n} u (\theta (x _ {i} ^ {*} \otimes y) (x _ {i})) \\ = \sum_ {i = 1} ^ {n} \theta (x _ {i} ^ {*} \otimes y) (u (x _ {i})) = \Big (\sum_ {i = 1} ^ {n} \langle u (x _ {i}), x _ {i} ^ {*} \rangle \Big) y. \end{array}
\]

Consequently, u belongs to \(A_{M}\) , and M is balanced.

COROLLARY 1. --- A free module is balanced.

This is clear if the module is zero, and a nonzero free module is generating (VIII, p.~81, Example 2).

COROLLARY 2. --- Let A be a ring, and let n be an integer \(\geqslant 0\) . The center of \(\mathbf{M}_{n}(\mathrm{A})\) consists of the scalar matrices with entries in the center of A. We view \(A^{n}\) as a left \(\mathbf{M}_{n}(\mathrm{A})\) -module (II, §10, No.~7, p.~349). The endomorphisms of this module are the mappings \(x \mapsto xa\) , where a runs through A.

Let M be the right A-module \(A_{d}^{n}\) . It is balanced by Corollary 1. Consequently, the centers of \(A_{M}\) , \(A_{M}^{\prime}\) , and \(A_{M}^{\prime\prime}\) coincide. Corollary 2 then follows from the fact that \(A_{M}\) is identified with A and \(A_{M}^{\prime}\) with \(\mathbf{M}_{n}(\mathbf{A})\) .

Remark. --- Let M be an A-module. If the \(A_{M}\) -module M is generating, then we have \(A_{M} = A_{M}^{\prime\prime}\) by Theorem 2 applied to the \(A_{M}\) -module M, so that M is a balanced A-module.

COROLLARY 3. --- Every finitely generated projective module over a commutative ring is balanced.

Indeed, a finitely generated projective module M is a generating \(A_{M}\) -module by Proposition 3 of VIII, p.~82. The corollary therefore follows from the remark above.

COROLLARY 4. --- Every finitely generated module over a principal ideal domain is balanced.

Indeed, a finitely generated module M over a principal ideal domain A is a generating \(A_{M}\) -module (VIII, p.~81, Example 4).

COROLLARY 5. --- Let K be a commutative field, V a finite-dimensional vector space over K, and u and v endomorphisms of V. The following properties are equivalent:

\begin{enumerate}
\def\labelenumi{(\roman{enumi})}
\item
  There exists a polynomial \(\mathrm{P}\) in \(\mathrm{K}[\mathrm{X}]\) such that \(v = \mathrm{P}(u)\) .
\item
  The endomorphism v commutes with every endomorphism of V that commutes with u.
\end{enumerate}

Take A to be the ring K{[}X{]} and M to be the K{[}X{]}-module deduced from V and u (VII, §5, No.~1, p.~28). Assertion (i) means that \(v \in K[X]_{M}\) and (ii) that \(v \in K[X]_{M}'\) . Corollary 5 is therefore a particular case of Corollary 4.

PROPOSITION 4. --- A semisimple module with finitely generated countermodule is balanced.

This follows from Example 3 of VIII, p.~81 and the remark.

COROLLARY 1. --- Let \((\mathrm{S}_{i})_{i\in\mathrm{I}}\) be a finite family of pairwise nonisomorphic simple A-modules. For \(i\in I\) , denote by \(D_{i}\) the opposite ring of the field of endomorphisms of \(S_{i}\) . Suppose that for every \(i\in I\) , the \(D_{i}\) -vector space \(S_{i}\) is finite-dimensional. Then the mapping \(a \mapsto (a_{\mathrm{S}_i})_{i \in \mathrm{I}}\) from A to \(\prod_{i \in \mathrm{I}} \mathrm{End}_{\mathrm{D}_i}(\mathrm{S}_i)\) is surjective.

Consider the A-module \(M = \prod_{i \in I} S_i\) . Since I is finite, we also have \(M = \bigoplus_{i \in I} S_i\) , and the image of \(S_i\) in M is the isotypical component of M of type \(S_i\) . Consequently, the endomorphisms of the A-module M are the mappings \((s_i) \mapsto (s_i d_i)\) , where \((d_i)_{i \in I}\) runs through \(\prod_{i \in I} D_i\) (VIII, p.~66, Proposition 5). Since I is finite and for every \(i \in I\) , the right \(D_i\) -vector space \(S_i\) is finite-dimensional, the countermodule of M is finitely generated. By Proposition 4, the A-module M is balanced. Now, the bicommutant of the A-module M consists of elements of \(\text{End}_{\mathbf{Z}}(\text{M})\) of the form \(\prod_{i \in I} u_i\) , where \((u_i) \in \prod_{i \in I} \text{End}_{\text{D}_i}(\text{S}_i)\) (VIII, p.~78, Proposition 1) because \(\text{End}_{\text{D}_i}(\text{S}_i)\) is the bicommutant of \(S_i\) . The corollary follows from this.

This corollary applies, in particular, when A is an algebra over a commutative field K and each \(S_{i}\) is a simple A-module that is finite-dimensional as a K-vector space: indeed, \(D_{i}\) then contains the homotheties \(\alpha_{S_{i}}\) , where \(\alpha\) runs through K, and \(S_{i}\) is finite-dimensional over \(D_{i}\) because it is so over K.

COROLLARY 2 (Burnside's theorem). --- Let A be an algebra over an algebraically closed commutative field K, and let S be a simple A-module that is finite-dimensional as a K-vector space. Then the mapping \(a \mapsto a_{\mathrm{S}}\) from A to \(\operatorname{End}_{\mathrm{K}}(\mathrm{S})\) is surjective.

Indeed, the field of endomorphisms of the A-module S consists of the homotheties \(\alpha_{S}\) with \(\alpha \in K\) (VIII, p.~47, Theorem 1). We then apply Corollary 1 to the simple A-module S.

